\section{When later disclosures change the historical query}
\label{sec:apple-pronouncements}

A later filing can reveal more about an earlier period without identifying a
different economic history. For example, a new segment-expense disclosure may
make a previously hidden cost visible, while the transactions that generated
that cost remain the same. To use that information correctly, a model must
retain the original report, the later presentation and the date at which the
additional detail became available. Apple's accounting-pronouncement discussion
shows why those distinctions are needed before the later figures even exist.

Annual and interim MD\&A describe planned changes to tax and segment disclosures;
the interim adds expense disaggregation and a subsequent effective-date
clarification \citep[Recent Accounting Pronouncements, printed p.~26]{apple2024}
\citep[Recent Accounting Pronouncements, printed p.~17]{apple2025q1}.
The following table preserves Apple's stated adoption plans. The quarters are
fiscal, and the descriptions are issuer reports of the updates, rather than an
independent determination of their operative requirements.
\begin{center}
\small
\begin{tabular}{p{0.19\textwidth}p{0.40\textwidth}p{0.25\textwidth}}
\toprule
Update & Reported disclosure change & Apple's stated plan\\\midrule
ASU 2023-07 & Segment expenses and information about the chief operating
decision maker & Fiscal 2025 Q4, retrospective\\
ASU 2023-09 & Income-tax reconciliation and taxes-paid detail & Fiscal 2026 Q4,
prospective\\
ASU 2024-03, clarified by 2025-01 & Expense components and selling-expense
amount and definition & Fiscal 2028 Q4, prospective\\\bottomrule
\end{tabular}
\end{center}

The tax discussion says the rate reconciliation will include additional
information and detail for items meeting a quantitative threshold. Taxes paid
will be disaggregated into federal, state and foreign taxes, with further detail
for significant individual jurisdictions. It gives neither the threshold value
nor the future jurisdictional amounts. Those cash-tax populations also remain
distinct from tax expense, payable and deferred balances; their reconciliation
is developed in Section~\ref{sec:apple-tax-settlement}.
An announced disclosure change cannot be used to allocate an old tax total
among jurisdictions by assumption.

The segment discussion concerns expenses that are significant and regularly
provided to the chief operating decision maker, together with that person's
title, position and use of segment profit or loss in assessing performance and
allocating resources. The future disclosure may answer questions about the
information used by management. It does not identify the regional allocation
of the existing corporate expense pool or change the current absence of
segment asset information. The previously examined retrospective-adoption
sentence remains a dated plan, not a later recast observation.

The interim expense passage reports ASU 2024-03's issuance in November 2024
and ASU 2025-01's clarification in January 2025. It describes purchases of
inventory, employee compensation, depreciation and intangible-asset amortization,
as applicable, within certain expense captions, a qualitative description of
the remaining amounts, and both the amount and the company's definition of
selling expenses. This describes disclosure populations; it does not provide
their values or make selling expenses identical to the entire selling, general
and administrative line. Purchases included in an expense disclosure likewise
cannot be identified with cash paid to suppliers without the relevant timing
and population facts.

These plans require separate historical query views. If a later filing presents
an earlier segment period retrospectively, an original-information query retains
the earlier disclosure, while a later-information query may use the checked
recast presentation. The original must not be overwritten. A comparison also
needs to establish the affected periods, definitions, units and basis before
matching the two versions. The words ``retrospective transition'' alone supply
neither a recast amount nor proof that a particular historical component will
be available. For the prospective tax and expense plans, these passages supply
no backward allocation of missing historical components. Exact transition
requirements need the operative updates and any applicable amendments.

The January clarification illustrates the timing problem directly. It occurs
after the interim period ended on 28 December 2024 but is described in the
interim report. An event can therefore belong to the information in a report
without occurring inside its reporting period. A historical forecast cutoff
needs evidence of availability and the exact retained version, not merely the
period-end or the update's issuance month. Similarly, an announced fiscal
adoption quarter is not an observed adoption event; the later filing must be
checked before the model records the plan as realised.

\subsection{Unchanged critical estimates do not mean constant amounts}

Annual Critical Accounting Estimates says management makes judgments,
assumptions and estimates affecting reported amounts, using historical
experience and other assumptions it considers reasonable to judge carrying
values. Its uncertain-tax discussion says reserves can be adjusted as facts
and circumstances change, and that final outcomes may differ from those
reserves. Its legal-contingency policy is the incurred-loss and estimability
condition examined above \citep[Critical Accounting Estimates, printed p.~26]{apple2024}.
These are descriptions of judgment and policy, not a disclosed numerical
estimation model or an assurance that the estimates will prove accurate.

The interim references both annual and interim Note 1 and reports no material
changes to critical accounting estimates since the annual filing
\citep[Critical Accounting Estimates, printed p.~17]{apple2025q1}.
That statement preserves the annual dependency and its materiality qualifier.
It cannot make every carrying amount, reserve or underlying assumption exactly
constant, and it supplies no numerical tolerance for a comparison. A model
should retain the reported policy and qualified continuity assertion alongside
the actual dated balances, then request the missing estimate or transition
facts when the query needs them. The present evidence supports those separate
records; it does not establish current compliance, realised adoption or the
historical expense components that future disclosures may reveal.
